% Kapitel 1 --- Einleitung
% Budget: 1,5 Seiten (von 13 ohne Literatur und Anhang)
% Quelle im Vault: 10 Thema/Thema 29 -- Überblick, 10 Thema/3.1 Konsistenz

\section{Einleitung}
\label{sec:einleitung}

\subsection{Problem}

Studieninformationen stehen nicht an einer Stelle. Dieselbe Angabe --- eine Bewerbungsfrist,
eine Zulassungsvoraussetzung, die Zahl der Leistungspunkte eines Moduls --- erscheint in der
Prüfungsordnung, auf der Studiengangsseite und in den Informationsschriften der zentralen
Studienberatung. Die Dokumente entstehen zu verschiedenen Zeiten, werden von verschiedenen
Stellen gepflegt und nur selten gemeinsam fortgeschrieben. Für die Leserin ist das
unsichtbar: Jede Seite wirkt für sich vollständig und aktuell, und es gibt keine Stelle, die
die Angaben gegeneinander hält.

Ein Beispiel aus dem Korpus dieser Arbeit zeigt, wie unauffällig solche Abweichungen
entstehen.\footnote{Übersichtsseite der Mitteilungen der JLU 7.36.07 Nr.~11 und PDF der
1.~Änderungsfassung; die Adressen sind in Abschnitt~\ref{sec:daten} verzeichnet.} Die
Spezielle Ordnung des Masterstudiengangs Data Analytics liegt in vier Fassungen vor. Für die
Fassung des 1.~Änderungsbeschlusses nennt die Übersichtsseite der Mitteilungen den
Geltungsbeginn \enquote{gültig ab dem Wintersemester 2023/24}; die Titelseite derselben
Fassung formuliert \enquote{gilt ab dem Wintersemester 2024/25}. Beide Angaben stehen auf
Servern derselben Hochschule, beide sind zitierfähig, und beide können nicht zutreffen. Wer
die Übersichtsseite liest, hält die Änderung ein Jahr früher für verbindlich als jemand, der
die Ordnung selbst aufschlägt.

Die Ursache ist strukturell und nicht persönlich: Die Übersichtsseite fasst zusammen, die
Ordnung ist der verbindliche Text, und beide werden von Hand gepflegt. Solche Stellen sind
der Gegenstand dieser Arbeit. Sie fragt nicht, ob die Universität ihre Seiten gut pflegt,
sondern ob sich Abweichungen dieses Typs systematisch finden, zuverlässig bewerten und nach
ihrer Entstehung unterscheiden lassen.

\subsection{Relevanz}

Informationsqualität ist keine Geschmacksfrage. ISO/IEC 25012 führt \emph{Consistency} als
eigenes Qualitätsmerkmal: Daten sind konsistent, wenn sie frei von Widersprüchen und mit
anderen Daten kohärent sind \parencite{iso25012}. Ein Datenprodukt, das sich selbst
widerspricht, ist damit unabhängig davon mangelhaft, ob jemand den Widerspruch bemerkt.

Wie schnell solche Widersprüche entstehen, ist für wissenschaftliche Verweise gut
dokumentiert. \textcite{klein2014} finden Reference Rot in jedem fünften der von ihnen
untersuchten Artikel; \textcite{jones2016} berichten, dass drei von vier URI-Verweisen auf
Inhalte zeigen, die sich gegenüber dem Zitat verändert haben. Die Studien betreffen
Literaturverweise und nicht Studieninformationen, und ihre Zahlen lassen sich nicht
übertragen. Sie belegen aber den Mechanismus: Dokumente werden überarbeitet, verschoben und
ersetzt, während die Verweise auf sie stehen bleiben.

Für eine Hochschule kommt ein zweiter Punkt hinzu. Eine Studienentscheidung ist folgenreich
und wird in der Regel einmal getroffen: Wer sich auf eine überholte Fassung stützt, merkt es
nicht an der Seite, sondern frühestens im Verfahren. Der Aufwand, solche Stellen zu finden,
ist heute Handarbeit --- und er wächst mit der Zahl der Dokumente. Genau dort setzt die
Automatisierung an, die diese Arbeit untersucht.

\subsection{Beitrag dieser Arbeit}

Die Arbeit leistet drei Beiträge. Erstens ein Verfahren, das Angaben aus verschiedenen
Dokumenten extrahiert, normalisiert und nach Sachverhalten gruppiert, sodass Abweichungen
systematisch sichtbar werden statt zufällig aufzufallen. Zweitens eine Messung der
Zuverlässigkeit dieses Verfahrens: Gegen einen von Hand erstellten Goldstandard werden
Precision und Recall je Sachverhaltstyp sowie die Übereinstimmung zwischen Mensch und Modell
berichtet. Drittens eine Typologie der Abweichung, die echte Widersprüche von
Versionsartefakten, Scope-Unterschieden und Extraktionsfehlern trennt --- erst sie macht aus
einer Fundliste einen Befund.

\subsection{Aufbau}

Kapitel~\ref{sec:hintergrund} ordnet die Arbeit in Informationsqualität,
Widerspruchserkennung und Content Drift ein und grenzt sie von den Nachbardimensionen ab.
Kapitel~\ref{sec:frage} formuliert Forschungsfrage und Hypothesen samt
Widerlegungskriterien. Kapitel~\ref{sec:daten} beschreibt Korpus, Faktkatalog und
Erhebungsgrenzen, Kapitel~\ref{sec:methode} die Analyseeinheit, die Extraktion mit
Zitatpflicht, die Normalisierung, die Abweichungstypologie und das Evaluationsdesign.
Kapitel~\ref{sec:ergebnisse} berichtet die Messung, Kapitel~\ref{sec:diskussion} ordnet sie
ein und benennt die Grenzen, Kapitel~\ref{sec:fazit} fasst das Ergebnis zusammen.
